\bookchapter{Apêndice H: roteiro de versões}

\chapterintro{Este livro continua sendo atualizado. Abaixo estão as direções planejadas para as próximas versões, e você também pode nos contar no seu feedback o conteúdo que mais gostaria de ver.}

\textbf{v1.8 já disponível: novas Partes VII, VIII e IX — “Prática de Programação com IA”, “Prática de Desenvolvimento Android” e “Prática de Desenvolvimento iOS” (Capítulos 37-69, 33 capítulos no total)}\par

Da revolução de paradigma da programação com IA, aos cinco princípios de pensamento, engenharia de contexto e de requisitos, panorama de ferramentas e modelos de 2026, prática com APIs, engenharia de diálogo, três projetos completos, garantia de qualidade e caminho de aprimoramento — levando você de “saber usar a IA” a “comandar a IA a escrever código”; até o ciclo completo de desenvolvimento nas duas grandes plataformas, Android e iOS: configuração do ambiente (Android Studio / Xcode), curso rápido das linguagens (Kotlin / Swift), frameworks de interface (Jetpack Compose / SwiftUI), dados e rede, projetos práticos (controle financeiro + consulta de previsão do tempo), testes e publicação (Google Play / App Store) e manuais de consulta rápida — levando você de usuário digital, passo a passo, a desenvolvedor de apps.

\begin{tcbinfo}{Planos para a próxima versão (v1.9)}
\begin{itemize}
  \item Celular e dispositivos vestíveis: guia de integração profunda com Apple Watch, smartwatches da Huawei e anéis inteligentes (Oura/Galaxy Ring)
  \item Viagens e celular: eSIM na prática, planos de dados internacionais, mapas offline, pagamentos no exterior e comparação de apps de tradução
  \item Celular e saúde: apps de exercícios, monitoramento do sono, frequência cardíaca e oxigenação do sangue, saúde feminina e saúde mental
  \item Criação de vídeos curtos com IA: tutoriais práticos de CapCut IA, Kling IA, avatares digitais e dublagem com IA
  \item Introdução a minisséries verticais: estrutura de roteiro, técnicas de gravação e políticas das plataformas
\end{itemize}

\end{tcbinfo}

\begin{tcbwarn}{Nota sobre a atualidade das informações}
A tecnologia de celulares evolui muito rápido, e este livro foi escrito com base na situação tecnológica e de mercado de meados de 2026. Recursos de IA, modelos de chips, versões de sistema e faixas de preço podem já ter mudado quando você estiver lendo. Em passos de operação específicos, siga a versão real do sistema do seu aparelho; em decisões de compra, consulte as análises e os comentários de usuários mais recentes.

\end{tcbwarn}
